\section{Agent Behavior Analysis}
\label{sec:behavior}

\paragraph{Intended behavior.}
The agent should preserve individual rationality, concede as delay becomes more
costly, escape repeated paths, and separate evidence by role. In persuasion it
should recommend high quality, pool low quality only within receiver obedience,
and distribute pooled signals smoothly across the horizon.

\paragraph{Executable alignment mechanisms.}
Four mechanisms connect intent to action. First, the policy reads only the
filtered game state and updates beliefs only from visible outcomes. Second,
offer and response validators enforce feasibility, budget balance, and legal
schemas. Third, cycle counters and terminal rules prevent strategically legal
but non-terminating behavior. Fourth, assignment-level audits apply separate
role and family exposure limits. Deterministic pacing in Eq.~\ref{eq:pacing}
makes every seller signal reproducible from the observed history; it controls
aggregate frequency rather than proving dynamic receiver obedience.

\paragraph{Observed alignment.}
V33 completed 150 bargaining games with positive sums for Alice and Bob, a sign
consistent with but not attributable to the two branch designs. In one live
unknown-horizon game, Bob ($\delta_B=0.95$) received 36.5079\% in round one.
The inherited branch rejected against a 43.32\% continuation threshold; the
recovery override accepted because the offer exceeded $f_B=33\%$, producing an
immediate agreement. The same game lost 4.98 rating points, illustrating that
agreement-aligned behavior and the competition metric can diverge. Persuasion
buyers and sellers were both positive, while offline grid tests only establish
that the computed marginal target does not exceed Eq.~\ref{eq:obedience}; they
do not establish live buyer posteriors. Zero action fallbacks rule out that
recorded execution-failure channel, not other policy or measurement errors.

\paragraph{Deviations and recovery.}
Alignment is not equivalent to rating success. Both negotiation roles lost,
although offers remained individually rational and contract-valid. The family
guard stopped exposure after 20 games. Earlier failures sharpen this
distinction: V3's 99-round loop was legal but behaviorally misaligned with
termination, while V11's 53-game return from a 24-game request was strategically
unchanged but operationally unsafe. These cases motivated cycle tests and the
one-assignment controller. V33 shows the remaining limitation: safeguards can
bound a bad trajectory, but cannot guarantee that a historically successful
policy remains profitable against a changing population.
